% =============================================================================
%  Notas de clase — Sesión 21: Archivos de texto
%  Programación Avanzada — Universidad Panamericana
%  Compilar: pdflatex sesion21_archivos_texto.tex
% =============================================================================
\documentclass[12pt, oneside]{article}
\usepackage{progavanzada}
\usepackage{array}

\begin{document}

\portada{Sesión 21}{Archivos de texto}{2026}

\tableofcontents
\newpage

% =============================================================================
\section{1970: cuando la computadora aprendió a tener cajones}
% =============================================================================

En 1970, en los laboratorios Xerox PARC en Palo Alto, un equipo de investigadores
liderado por Alan Kay se enfrentaba a un problema que ninguna ecuación podía
resolver: las computadoras eran intimidantes.  Los usuarios tenían que memorizar
comandos crípticos, navegar sistemas de archivos sin ningún contexto visual y
entender abstracciones que no tenían nada que ver con su vida cotidiana.

La solución fue radical: hacer que la computadora pareciera un escritorio de
oficina.  Un escritorio con papeles encima, cajones donde guardar documentos,
carpetas para organizar los cajones y una papelera de reciclaje en la esquina.
No fue una decisión de diseño gráfico.  Fue una decisión pedagógica deliberada:
si la gente ya sabe usar un escritorio, puede aprender a usar una computadora.

El nombre de los conceptos no es traducción: \textit{desktop}, \textit{folder},
\textit{file} y \textit{trash} fueron elegidos precisamente porque ya existían
en el mundo físico.  Hoy, en esta sesión, vas a entender por qué esa metáfora
es tan precisa: porque la diferencia entre la memoria de tu programa y un
archivo en disco es exactamente la diferencia entre poner un papel sobre el
escritorio y guardarlo en un cajón.

\begin{tecnico}[El escritorio digital: Xerox PARC]
  Alan Kay y su equipo desarrollaron entre 1970 y 1973 el sistema
  \textbf{Smalltalk}, el primero en combinar ventanas, iconos, ratón y menús
  —lo que hoy conocemos como interfaz gráfica.  Apple llevó esas ideas al
  mercado masivo con el Macintosh en 1984; Microsoft con Windows en 1985.
  \par\smallskip
  Antes de Smalltalk, los archivos se llamaban \textit{datasets} y se
  referenciaban por códigos internos.  El término \textit{file} fue elegido
  deliberadamente para evocar los archiveros físicos de oficina.
  \par\smallskip
  \textit{Referencia:} Kay, A. \& Goldberg, A. (1977). ``Personal Dynamic
  Media.'' \textit{IEEE Computer}, 10(3), 31--41.
\end{tecnico}

% =============================================================================
\section{Memoria vs.\ archivos: el escritorio y los cajones}
% =============================================================================

Cuando tu programa declara una variable, la computadora le asigna espacio en
la \textbf{memoria RAM}.  Puedes leer ese valor, modificarlo, usarlo en
cálculos.  Es rápido y cómodo, como tener un papel justo frente a ti sobre
el escritorio.

El problema es que el escritorio se limpia solo cuando el programa termina.
Todas las variables, todos los arreglos, todos los cálculos que hiciste
desaparecen.  La próxima vez que ejecutes el programa, empiezas de cero.

Los \textbf{archivos} resuelven esto.  Un archivo vive en el disco duro (o en
un SSD), no en RAM.  Cuando el programa termina, el archivo sigue ahí.  Cuando
apagas la computadora, el archivo sigue ahí.  Es como guardar un papel en un
cajón: el cajón no se vacía solo.

\begin{center}
\begin{tabular}{lll}
  \toprule
  \textbf{Característica}   & \textbf{Memoria (RAM)} & \textbf{Archivo (disco)} \\
  \midrule
  Acceso                    & Instantáneo            & Más lento \\
  Persiste al cerrar        & No                     & Sí \\
  Persiste al apagar        & No                     & Sí \\
  Capacidad típica          & Gigabytes              & Terabytes \\
  Metáfora                  & Escritorio             & Cajón \\
  \bottomrule
\end{tabular}
\end{center}

\begin{recuerda}
  Un programa puede guardar resultados en un archivo y otro programa puede
  leerlos después.  Esa es la base de todo sistema de persistencia: bases
  de datos, configuraciones, imágenes, música, documentos.  Todo es, al
  final, un archivo.
\end{recuerda}

% =============================================================================
\section{¿Qué es un archivo, exactamente?}
% =============================================================================

Un archivo es una \textbf{secuencia de bytes} almacenada en disco.  Nada más.
No tiene estructura mágica, no sabe qué contiene, no se protege solo.  Lo que
hace que un archivo sea texto, imagen o música es únicamente cómo decidimos
interpretar esos bytes.

\subsection{Verlo con tus propios ojos}

Antes de escribir una sola línea de C++, crea un archivo de texto a mano:

\begin{description}
  \item[Windows] Clic derecho en el escritorio $\to$ Nuevo $\to$ Documento de texto.
        Escribe algunas letras (por ejemplo, \texttt{Hola}) y guarda.
  \item[Linux / Mac] En la terminal: \texttt{touch prueba.txt}, luego ábrelo
        con cualquier editor y escribe lo mismo.
\end{description}

Ahora ábrelo con un \textbf{editor hexadecimal en línea} (por ejemplo,
\texttt{hexed.it}).  Lo que ves es el contenido real del archivo: cada
carácter representado por su valor en hexadecimal.

La \texttt{H} de ``Hola'' es \texttt{0x48}, la \texttt{o} es \texttt{0x6F},
la \texttt{l} es \texttt{0x6C}, la \texttt{a} es \texttt{0x61}.  Son los
valores de la tabla ASCII que ya conoces.

\begin{center}
\begin{tabular}{cccc}
  \toprule
  \textbf{Carácter} & \textbf{Decimal} & \textbf{Hexadecimal} & \textbf{Binario} \\
  \midrule
  \texttt{H} & 72  & \texttt{0x48} & \texttt{01001000} \\
  \texttt{o} & 111 & \texttt{0x6F} & \texttt{01101111} \\
  \texttt{l} & 108 & \texttt{0x6C} & \texttt{01101100} \\
  \texttt{a} & 97  & \texttt{0x61} & \texttt{01100001} \\
  \bottomrule
\end{tabular}
\end{center}

Fíjate en el tamaño del archivo: 4~bytes para 4~caracteres.  Agrega una
letra más y el archivo crece exactamente 1~byte.

\begin{tecnico}[El estándar lo garantiza: \texttt{sizeof(char) == 1}]
  En C y C++, el estándar garantiza que \cpp{sizeof(char) == 1} en cualquier
  plataforma, compilador o arquitectura.  Un \cpp{char} siempre ocupa
  exactamente 1~byte.  Por eso un archivo de texto con $n$ caracteres ocupa
  exactamente $n$~bytes (más los saltos de línea).
  \par\smallskip
  Curiosidad histórica: el primer disco duro comercial, el IBM RAMAC~350
  (1956), almacenaba 5~MB, pesaba una tonelada y se rentaba por
  USD~3\,200 al mes.  Una microSD actual de 1~TB cuesta menos de USD~100
  y cabe en la uña del pulgar.  El mismo concepto de ``archivo en disco'',
  200\,000 veces más capacidad.
  \par\smallskip
  \textit{Referencia:} IBM Archives. ``IBM 350 disk storage unit.''
  \texttt{ibm.com/ibm/history/exhibits/storage}
\end{tecnico}

\subsection{Saltos de línea}

Un salto de línea también es un carácter (o dos).  Aquí hay una diferencia
importante entre sistemas operativos:

\begin{center}
\begin{tabular}{lll}
  \toprule
  \textbf{Sistema} & \textbf{Secuencia} & \textbf{Bytes} \\
  \midrule
  Linux / Mac   & \texttt{\textbackslash n}              & \texttt{0x0A} (1 byte) \\
  Windows       & \texttt{\textbackslash r\textbackslash n} & \texttt{0x0D 0x0A} (2 bytes) \\
  \bottomrule
\end{tabular}
\end{center}

\begin{advertencia}
  Si creas un archivo de texto en Windows y lo abres en Linux (o viceversa),
  puedes ver caracteres extraños al final de cada línea.  Esto es causa de
  muchos bugs cuando se comparten archivos entre sistemas.  La biblioteca
  \cpp{<fstream>} en modo texto maneja la conversión automáticamente en la
  mayoría de los casos, pero es bueno saber que el problema existe.
\end{advertencia}

% =============================================================================
\section{La biblioteca \texttt{<fstream>}}
% =============================================================================

Para leer y escribir archivos en C++ se usa la biblioteca \textbf{\cpp{<fstream>}}
(\textit{file stream}).  Provee tres clases:

\begin{center}
\begin{tabular}{lll}
  \toprule
  \textbf{Clase}   & \textbf{Uso}             & \textbf{Analogía}          \\
  \midrule
  \cpp{ifstream}   & Leer de un archivo       & Sacar un papel del cajón   \\
  \cpp{ofstream}   & Escribir en un archivo   & Meter un papel al cajón    \\
  \cpp{fstream}    & Leer y escribir          & Cajón de doble apertura    \\
  \bottomrule
\end{tabular}
\end{center}

La idea central es que un archivo se comporta igual que \cpp{cin} y \cpp{cout}:
usas los mismos operadores \cpp{>>} y \cpp{<<}, solo que aplicados a un objeto
que representa un archivo en disco en lugar de la terminal.

\begin{consejo}
  Si ya sabes usar \cpp{cin} y \cpp{cout}, ya sabes casi todo lo necesario
  para leer y escribir archivos.  \cpp{cin} y \cpp{cout} son también
  \textbf{streams}: el primero conectado al teclado, el segundo a la pantalla.
  \cpp{ifstream} y \cpp{ofstream} son lo mismo pero conectados a un archivo
  en disco.  El operador cambia de destino, no de sintaxis.
\end{consejo}

% =============================================================================
\section{Leyendo un archivo línea por línea}
% =============================================================================

El primer ejemplo lee un archivo de texto y muestra su contenido en pantalla.

\textbf{canciones.txt}:
\begin{verbatim}
Walking on Sunshine
Happy
Color Esperanza
\end{verbatim}

\textbf{archivos1.cpp}:
\begin{verbatim}
#include <iostream>
#include <fstream>
#include <string>
using namespace std;

int main(){
    ifstream ifile("canciones.txt");
    string cancion;

    while (getline(ifile, cancion)){
        cout << cancion << endl;
    }
}
\end{verbatim}

Desglose:

\begin{description}
  \item[\cpp{ifstream ifile("canciones.txt");}]
    Crea un objeto de tipo \textbf{\cpp{ifstream}} y lo conecta al archivo
    \texttt{canciones.txt}.  Si el archivo no existe, el objeto queda en
    estado de error pero el programa no truena por defecto.

  \item[\cpp{while (getline(ifile, cancion))}]
    Lee una línea completa del archivo y la almacena en \cpp{cancion}.
    El salto de línea se consume pero \textit{no} se incluye en la cadena.
    El \cpp{while} se detiene cuando no quedan más líneas: \cpp{getline}
    devuelve el propio stream, y al llegar al final del archivo el stream
    se evalúa como \cpp{false}.

  \item[\cpp{cout << cancion << endl;}]
    Muestra la línea en pantalla.  El \cpp{endl} agrega el salto que
    \cpp{getline} ya consumió del archivo.
\end{description}

\begin{advertencia}
  El archivo debe estar en el mismo directorio que el \textbf{ejecutable}
  (no necesariamente donde está el \texttt{.cpp}).  Si el archivo no se
  encuentra, el \cpp{while} no itera nunca y el programa termina en
  silencio.  Para detectarlo:
  \begin{verbatim}
ifstream ifile("canciones.txt");
if (!ifile){
    cout << "Error: no se pudo abrir el archivo" << endl;
    return 1;
}
  \end{verbatim}
\end{advertencia}

% =============================================================================
\section{Escribiendo en un archivo}
% =============================================================================

Ahora al revés: en lugar de leer desde un archivo, escribimos resultados en uno.

\textbf{archivos2.cpp}:
\begin{verbatim}
#include <iostream>
#include <fstream>
#include <string>
using namespace std;

int main(){
    ofstream ofile("resultados.txt");
    int x = 10;
    int y = 20;

    ofile << x+y << " ";
    ofile << y-x << " ";
    ofile << (x+y)/2 << endl;

    ofile.close();
}
\end{verbatim}

Después de ejecutar este programa, \texttt{resultados.txt} contendrá:
\begin{verbatim}
30 10 15
\end{verbatim}

El operador \cpp{<<} es exactamente el mismo que usas con \cpp{cout}.  La
única diferencia es que el destino es un archivo en lugar de la pantalla.

\subsection{¿Por qué llamar a \texttt{close()}?}

El sistema operativo usa un \textbf{buffer}: acumula los datos antes de
escribirlos al disco para hacer menos operaciones físicas, que son lentas.
Al llamar \cpp{close()}, le dices que vacíe el buffer y garantice que todo
llegó al disco.  Si el programa termina sin cerrar, el buffer se vacía
automáticamente, pero llamar \cpp{close()} explícitamente es buena práctica.

\begin{recuerda}
  Si \texttt{resultados.txt} ya existía, \cpp{ofstream} lo
  \textbf{sobreescribe} desde el principio.  Si quieres \textit{agregar}
  contenido al final sin borrar lo anterior, usa el modo \cpp{fstream::app}
  (\textit{append}):
  \begin{verbatim}
ofstream ofile("resultados.txt", fstream::app);
  \end{verbatim}
\end{recuerda}

% =============================================================================
\section{Leyendo valores numéricos}
% =============================================================================

El programa anterior escribió tres números en \texttt{resultados.txt}.
Ahora los leemos con otro programa.

\textbf{archivos3.cpp}:
\begin{verbatim}
#include <iostream>
#include <fstream>
#include <string>
using namespace std;

int main(){
    ifstream fresultados("resultados.txt");
    int x, y, z;
    fresultados >> x >> y >> z;
    cout << "La suma de " << x << " " << y << " y " << z;
    cout << " es " << x+y+z << endl;
}
\end{verbatim}

El operador \cpp{>>} sobre un \cpp{ifstream} funciona igual que sobre
\cpp{cin}: salta espacios y saltos de línea, lee el siguiente token y lo
convierte al tipo indicado.

\begin{ejemplo}[El encadenamiento funciona igual que con \texttt{cin}]
  Las tres lecturas se pueden escribir separadas o encadenadas:
  \begin{verbatim}
fresultados >> x;
fresultados >> y;
fresultados >> z;
// equivale a:
fresultados >> x >> y >> z;
  \end{verbatim}
  Cada \cpp{>>} devuelve el stream, lo que permite el encadenamiento.
  Con \cpp{cin} lo haces todo el tiempo sin pensarlo.
\end{ejemplo}

\begin{consejo}
  Nota el flujo completo de los tres ejemplos:
  \textbf{archivos2} \textit{escribe} \texttt{resultados.txt};
  \textbf{archivos3} lo \textit{lee} desde otro programa.  Ese es el poder
  de los archivos: la comunicación entre programas, sesiones y usuarios
  distintos.
\end{consejo}

% =============================================================================
\section{Actividad: agenda de cumpleaños}
% =============================================================================

Vamos a construir una pequeña base de datos colaborativa.  Cada persona
genera un archivo con su nombre y fecha de cumpleaños, lo comparte con el
grupo, y entre todos escriben programas para consultarla.

\subsection{Formato del archivo}

Cada archivo personal tiene cuatro líneas, en este orden:

\begin{verbatim}
nombre
apellidos
día
mes
\end{verbatim}

Por ejemplo, el archivo de Hannah Fabián (cumpleaños el 29 de febrero)
sería:

\begin{verbatim}
Hannah
Fabián
29
2
\end{verbatim}

\subsection{Parte A: genera tu archivo (10 min)}

Escribe un programa que genere tu propio archivo usando lo que aprendiste
en \textbf{archivos2.cpp}.  Nómbralo con tu nombre en minúsculas y sin
espacios, por ejemplo \texttt{hannah.txt}.  Una vez generado, súbelo a la
carpeta de Google Drive indicada por el profesor.

\subsection{Parte B: las consultas (40 min)}

Descarga todos los archivos de tus compañeros y colócalos en una carpeta.
Luego resuelve los siguientes problemas.  Para cada uno escribe un programa
separado.

\begin{enumerate}

  \item \textbf{Unir todos los datos en un solo archivo.}\\
        Lee cada archivo individual y escribe toda la información en un
        archivo llamado \texttt{agenda.txt}, una persona por línea con el
        formato:\\[4pt]
        \texttt{nombre apellidos día mes}\\[4pt]
        Puedes apoyarte en un archivo de texto auxiliar que liste los
        nombres de todos los archivos individuales, uno por línea.

  \item \textbf{¿Quién cumple años en un mes dado?}\\
        Lee \texttt{agenda.txt}, pide un número de mes al usuario (1~=~enero,
        12~=~diciembre) y muestra el nombre completo de todas las personas
        que cumplen años ese mes.

  \item \textbf{¿Cuál es el siguiente cumpleaños a partir de hoy?}\\
        Hoy es 6 de mayo.  ¿De quién es el próximo cumpleaños?  Si ningún
        compañero cumple años después del 6 de mayo, el siguiente es el
        primero del año que viene.  Si varias personas comparten la fecha,
        muéstralas todas.

  \item \textbf{¿Qué par de compañeros tiene los cumpleaños más cercanos?}\\
        Encuentra el par de personas cuyas fechas de cumpleaños están
        separadas por el menor número de días (considera el año como circular:
        el 31 de diciembre y el 1 de enero están separados por 1 día).

\end{enumerate}

% =============================================================================
\section{Ejercicios}
% =============================================================================

\begin{enumerate}

  \item Escribe un programa que pida al usuario tres calificaciones, las
        guarde en un archivo llamado \texttt{calificaciones.txt} (una por
        línea) y luego lea ese mismo archivo y muestre el promedio.

  \item El programa del ejercicio anterior sobreescribe el archivo cada
        vez que se ejecuta.  Modifícalo para que \textit{agregue} las
        calificaciones al final sin borrar las anteriores.  Después de
        agregar, lee todo el archivo y calcula el promedio de
        \textit{todas} las calificaciones almacenadas.

  \item Escribe un programa que lea \texttt{canciones.txt} y genere un
        archivo llamado \texttt{canciones\_numeradas.txt} donde cada
        línea tenga el número de línea seguido de la canción:
        \begin{verbatim}
1. Walking on Sunshine
2. Happy
3. Color Esperanza
        \end{verbatim}

  \item Escribe un programa que lea dos archivos de texto (cuyos nombres
        pide al usuario) y genere un tercer archivo que contenga el
        contenido del primero seguido del contenido del segundo.

  \item ¿Qué ocurre si intentas leer de un archivo que no existe con
        \cpp{ifstream}?  ¿Lanza una excepción?  ¿Da error de compilación?
        ¿Simplemente no lee nada?  Primero razona la respuesta, luego
        verifica experimentalmente.

\end{enumerate}

% =============================================================================
\section*{Resumen}
% =============================================================================

\begin{itemize}
  \item La \textbf{memoria RAM} es volátil: los datos desaparecen cuando
        el programa termina.  Los \textbf{archivos} en disco persisten.
  \item Un archivo es una secuencia de bytes.  En texto, cada carácter
        ocupa exactamente 1~byte (valor ASCII).
  \item \cpp{<fstream>} provee \cpp{ifstream} (leer), \cpp{ofstream}
        (escribir) y \cpp{fstream} (ambos).
  \item Los operadores \cpp{>>} y \cpp{<<} funcionan igual con archivos
        que con \cpp{cin} y \cpp{cout}: misma sintaxis, distinto destino.
  \item \cpp{getline(stream, cadena)} lee una línea completa; el
        \cpp{while} itera hasta el final del archivo.
  \item \cpp{ofstream} sobreescribe el archivo por defecto; usa
        \cpp{fstream::app} para agregar al final.
  \item Llama a \cpp{close()} al terminar de escribir para garantizar
        que el buffer se vacíe a disco.
  \item En Windows los saltos de línea son \texttt{\textbackslash r\textbackslash n}
        (2~bytes); en Linux/Mac son \texttt{\textbackslash n} (1~byte).
\end{itemize}

\end{document}
